%%%%%%%%%%%%%%%%%%%%%%%%%%%%%%%%%%%%%%%%%%%%%%%%%%%%%%%%%%%%%%%%%%%%%%%%%%%%%%
%
% 07_geometry.tex
%
%%%%%%%%%%%%%%%%%%%%%%%%%%%%%%%%%%%%%%%%%%%%%%%%%%%%%%%%%%%%%%%%%%%%%%%%%%%%%%

\begin{mosbox}{6. Semantic Geometry}

\BodyFont

The combinatorial structure of the simplicial complex and the local
constraints imposed by the cellular sheaf describe how information is
organized.

Semantic geometry introduces a complementary question:

\[
\boxed{
\text{How should concepts themselves be represented?}
}
\]

Within MOS, every concept is modeled probabilistically as an uncertainty
distribution rather than a single point in space.

This representation preserves both the estimated semantic location and the
associated confidence of that estimate.

\end{mosbox}

%%%%%%%%%%%%%%%%%%%%%%%%%%%%%%%%%%%%%%%%%%%%%%%%%%%%%%%%%%%%%%%%%%%%%%%%%%%%%%
%% GAUSSIAN REPRESENTATION
%%%%%%%%%%%%%%%%%%%%%%%%%%%%%%%%%%%%%%%%%%%%%%%%%%%%%%%%%%%%%%%%%%%%%%%%%%%%%%

\begin{eqbox}

\[
\boxed{
X_i
\sim
\mathcal N(\mu_i,\Sigma_i)
}
\]

\BodyFont

Each concept is represented by

\[
(\mu_i,\Sigma_i),
\]

where

\begin{itemize}

\item $\mu_i$ denotes the semantic embedding,

\item $\Sigma_i$ represents uncertainty,

\item higher confidence corresponds to smaller covariance.

\end{itemize}

\end{eqbox}

%%%%%%%%%%%%%%%%%%%%%%%%%%%%%%%%%%%%%%%%%%%%%%%%%%%%%%%%%%%%%%%%%%%%%%%%%%%%%%
%% SEMANTIC DIAGRAM
%%%%%%%%%%%%%%%%%%%%%%%%%%%%%%%%%%%%%%%%%%%%%%%%%%%%%%%%%%%%%%%%%%%%%%%%%%%%%%

\vspace{5mm}

\begin{center}

\begin{tikzpicture}[scale=1]

% Left Cloud
\node[cloud, cloud puffs=13, cloud puff arc=120, aspect=1.5, draw=MOSRed, line width=2pt, minimum width=3.5cm, minimum height=2.5cm, fill=MOSBlack] (c1) at (-4,0) {};
\fill[MOSRed] (-4,0) circle(3pt);
\node[white, below=1.5cm] at (-4,0) {$(\mu_1,\Sigma_1)$};

% Right Cloud
\node[cloud, cloud puffs=12, cloud puff arc=120, aspect=1.4, draw=MOSRed, line width=2pt, minimum width=3cm, minimum height=2cm, fill=MOSBlack] (c2) at (4,0) {};
\fill[MOSRed] (4,0) circle(3pt);
\node[white, below=1.5cm] at (4,0) {$(\mu_2,\Sigma_2)$};

% Wasserstein Distance
\draw[<->, MOSGold, line width=2.5pt] (c1.east) -- (c2.west) node[midway, above=3mm, MOSGold] {\Large $W_2$};

\end{tikzpicture}

\end{center}

%%%%%%%%%%%%%%%%%%%%%%%%%%%%%%%%%%%%%%%%%%%%%%%%%%%%%%%%%%%%%%%%%%%%%%%%%%%%%%
%% DISTANCE
%%%%%%%%%%%%%%%%%%%%%%%%%%%%%%%%%%%%%%%%%%%%%%%%%%%%%%%%%%%%%%%%%%%%%%%%%%%%%%

\begin{mosbox}{Semantic Distance}

\BodyFont

The similarity between two concepts is measured geometrically rather than
combinatorially.

MOS adopts the Wasserstein metric to compare probabilistic semantic
representations.

\[
\boxed{
W_2
\left(
X_1,
X_2
\right)
}
\]

The resulting distance reflects differences in both semantic position and
uncertainty.

Unlike Euclidean distance alone, this comparison incorporates the full
probabilistic structure of the concepts.

\end{mosbox}

%%%%%%%%%%%%%%%%%%%%%%%%%%%%%%%%%%%%%%%%%%%%%%%%%%%%%%%%%%%%%%%%%%%%%%%%%%%%%%
%% FUSION
%%%%%%%%%%%%%%%%%%%%%%%%%%%%%%%%%%%%%%%%%%%%%%%%%%%%%%%%%%%%%%%%%%%%%%%%%%%%%%

\begin{mosbox}{Precision-Weighted Fusion}

\BodyFont

When multiple compatible concepts are combined, MOS performs a
precision-weighted fusion.

Concepts with greater certainty contribute proportionally more to the
resulting estimate.

The fused representation is obtained through precision averaging.

\[
\boxed{
\mu
=
\left(
\sum_i
\Sigma_i^{-1}
\right)^{-1}
\left(
\sum_i
\Sigma_i^{-1}\mu_i
\right)
}
\]

The associated covariance decreases as mutually consistent evidence is
combined.

\end{mosbox}

%%%%%%%%%%%%%%%%%%%%%%%%%%%%%%%%%%%%%%%%%%%%%%%%%%%%%%%%%%%%%%%%%%%%%%%%%%%%%%
%% VISUAL FUSION
%%%%%%%%%%%%%%%%%%%%%%%%%%%%%%%%%%%%%%%%%%%%%%%%%%%%%%%%%%%%%%%%%%%%%%%%%%%%%%

\begin{center}

\begin{tikzpicture}[scale=1]

% Cloud 1 (Faint)
\node[cloud, cloud puffs=11, cloud puff arc=120, aspect=1.5, draw=MOSRed, line width=1pt, minimum width=3cm, minimum height=2cm, opacity=0.5] at (-4,0) {};

% Cloud 2 (Faint)
\node[cloud, cloud puffs=10, cloud puff arc=120, aspect=1.3, draw=MOSRed, line width=1pt, minimum width=2.5cm, minimum height=1.8cm, opacity=0.5] at (0,0) {};

\draw[arrow, opacity=0.5] (-2.2,0) -- (-1.2,0);
\draw[arrow, opacity=0.5] (1.2,0) -- (2.2,0);

% Fused Cloud (Strong, Gold)
\node[cloud, cloud puffs=14, cloud puff arc=120, aspect=1.6, draw=MOSGold, line width=3pt, minimum width=2cm, minimum height=1.5cm, fill=MOSBlack] (c3) at (4,0) {};
\fill[MOSGold] (4,0) circle(4pt);
\node[MOSGold, below=1.2cm] at (4,0) {\bfseries Fused Concept};

\end{tikzpicture}

\end{center}

%%%%%%%%%%%%%%%%%%%%%%%%%%%%%%%%%%%%%%%%%%%%%%%%%%%%%%%%%%%%%%%%%%%%%%%%%%%%%%
%% INTERPRETATION
%%%%%%%%%%%%%%%%%%%%%%%%%%%%%%%%%%%%%%%%%%%%%%%%%%%%%%%%%%%%%%%%%%%%%%%%%%%%%%

\begin{highlightbox}

\SectionTitle{Interpretation}

\BodyFont

Semantic geometry complements the topological organization of MOS.

\begin{center}

\begin{tabular}{ccl}

Topology
&
$\Longrightarrow$
&
Which concepts are related?

\\[3mm]

Geometry
&
$\Longrightarrow$
&
How similar are they?

\\[3mm]

Probability
&
$\Longrightarrow$
&
How certain is that relationship?

\end{tabular}

\end{center}

Together, these components provide a mathematical description of both the
structure and the uncertainty of conceptual memory.

\end{highlightbox}

%%%%%%%%%%%%%%%%%%%%%%%%%%%%%%%%%%%%%%%%%%%%%%%%%%%%%%%%%%%%%%%%%%%%%%%%%%%%%%
%% REMARK
%%%%%%%%%%%%%%%%%%%%%%%%%%%%%%%%%%%%%%%%%%%%%%%%%%%%%%%%%%%%%%%%%%%%%%%%%%%%%%

\begin{remarkbox}

\BodyFont

Semantic geometry does not replace the simplicial complex or the cellular
sheaf.

Instead, it enriches the existing topological framework by introducing a
continuous geometric layer capable of measuring similarity, uncertainty and
information fusion.

The topological and geometric components therefore play complementary
roles within the overall MOS architecture.

\end{remarkbox}